\abstract{
Despite recent progress of VLA foundation models, the disparity between laboratory conditions and real-world applications continues to impede their practical implementation.
To bridge this gap, we present \method, which advances \methodold through improvements in three functional domains.
(1) \textbf{\textit{Generalization}} across tasks and embodiments.
Compared to the previous version, we revamp the data processing pipeline and curate around 60,000 hours of data for pretraining, including 50,000 hours of robot trajectories spanning \robotnum robot configurations and 10,000 hours of egocentric human videos.
(2) \textbf{\textit{Expanded action space}} in addition to dual-arm hardware platforms.
In particular, our system accommodates degrees of freedom for the heads, waists, mobile bases, and dexterous hands, thereby empowering the robots to tackle more complex tasks in practical scenarios.
(3) \textbf{\textit{Predictive dynamics modeling}} for improved temporal reasoning.
Specifically, we formulate future prediction as a proxy task, facilitated by a video representation model for semantic priors and a depth estimation model for geometric cues.
Evaluations on the GM-100 benchmark, conducted in a generalist setting, validate the beneficial impact of these proposed modifications.
Furthermore, benefiting from the expanded pretraining data that covers whole-body degrees of freedom, LingBot-VLA-2.0 demonstrates strong cross-embodiment long-horizon mobile manipulation capability across the two robotic platforms.


}
